\documentclass[12pt,a4paper]{article}
\usepackage{graphicx}
\usepackage{amsmath}
\usepackage{hyperref}
\usepackage{geometry}
\geometry{margin=2.5cm}
\emergencystretch=3em

\title{\textbf{CDFD Runtime Paper 07: Autonomous Discovery, Hypothesis Triage, and Falsification}}
\author{Steve Bico Mujjabi, MD\\
\small Independent Researcher\\
\small Founder, Vura Labs\\
\small Kampala, Uganda\\
\small ORCID: \href{https://orcid.org/0009-0001-0556-5516}{0009-0001-0556-5516}}
\date{June 2026}

\begin{document}
\maketitle

\begin{abstract}
\noindent \textbf{Executive synthesis.} The discovery layer is a hypothesis
factory, not an oracle. It perturbs runtime states, extracts features, proposes
candidate regimes, and records falsification targets. This paper aligns the
discovery manuscripts with the actual local experiment surface in
\texttt{cdfd\_runtime/experiments}: OOL phase mapping, finance-climate coupling,
schema invention, vacuum hysteresis, knot stability, adaptive-surface
instability, frontier sweeps, and gigamarathon coverage.
\end{abstract}

\paragraph{Greek-symbol convention.}
Unless a manuscript states a tighter equation-local meaning, Greek symbols are local CDFD variables or coefficients, not universal constants. Here $\Phi$ denotes driving flux, $\Psi_s$ denotes the adaptive operating ratio, $\Omega$ denotes a domain or state space, and $\Lambda$ denotes the Life Number where used. The coefficients $\alpha$, $\beta$, and $\gamma$ denote accumulation or reinforcement, relaxation or decay, and diffusion, coupling, or redistribution. The symbols $\delta$ and $\epsilon$ denote perturbation or tolerance scales, while $\eta$, $\lambda$, $\mu$, $\nu$, $\rho$, $\sigma$, $\tau$, and $\phi$ denote local efficiency, rate, baseline, frequency, density or correlation, noise or variance, time-scale, and phase or field terms. Subscripts restrict any coefficient to the named pathway, tissue, domain, or experiment.


\begin{figure}[ht]
\centering
\fbox{\begin{minipage}{0.92\textwidth}
\centering
\textbf{Runtime evidence path}\\[0.4em]
Prior CDFD mathematics $\longrightarrow$ CDFL/runtime state $\longrightarrow$ bounded output $\longrightarrow$ falsification target.
\end{minipage}}
\caption{Conceptual schematic for the runtime papers. The engine translates the mathematical frame from the prior CDFD papers into executable CDFL/runtime diagnostics; candidate outputs remain tied to explicit tests before any stronger claim is made.}
\end{figure}

\section{Prior CDFD Basis}
This manuscript does not rederive the mathematical core of CDFD. It uses the
Mujjabi Adaptive Operating Ratio and the constrained-vacuum notation introduced
in Part I \cite{MujjabiI2026}, the torus-knot hierarchy and boundary-mode work
from the Part I sequence \cite{MujjabiVI2026,MujjabiIX2026}, the public balance
law developed in the Part I capstone \cite{MujjabiX2026}, the Part I transport
tests \cite{MujjabiXII2026}, and the Life Number and tri-regime biology bridge
from Part II \cite{MujjabiPartII2026}. The runtime papers therefore act as an
execution layer: they keep the mathematics connected to commands, outputs,
falsification tests, and domain-specific limits.


\section{Discovery Discipline}
Autonomous discovery is useful only when it remains falsifiable. The CDFD
The Runtime discovery layer does not turn a pattern into a law merely because
the engine generated it. A discovery candidate needs:
\begin{itemize}
    \item a script or CDFL model;
    \item saved output;
    \item parameter values;
    \item a pattern label;
    \item a falsification condition;
    \item a rerun command.
\end{itemize}

\section{Discovery Modules}
The discovery package contains feature extraction, hypothesis, generator,
knowledge, sensitivity, detector, validator, and causal-discovery modules. The
workflow is:
\begin{enumerate}
    \item generate or perturb a state;
    \item run bounded dynamics;
    \item extract features such as mean $\Psi_s$, variance, collapse count, and
    stability index;
    \item assign a candidate pattern label;
    \item write a report or output file;
    \item attempt falsification by changing parameters or null-coupling.
\end{enumerate}

\section{Selected Local Discoveries}
\subsection{Discovery A: OOL phase boundary}
\texttt{experiments/notebooks/discovery\_ool\_phase.py} sweeps a two-dimensional
prebiotic parameter space. Its selected HDF5 output contains grid points near
$\Lambda=1$. This is useful as a cross-Part-II phase-boundary bridge and as a
test target for the Life Number. It is not empirical proof of abiogenesis.

\subsection{Discovery B: finance-climate coupling}
\texttt{experiments/notebooks/discovery\_finance\_climate.py} models a coupled
finance-climate toy system and then runs causal discovery over the time series.
The report frames the edge from financial $\Psi$ volatility to climate
constraint as a candidate causal result. The necessary null test is explicit:
set coupling to zero and verify that the edge disappears.

\subsection{Discovery C: schema invention}
The schema-invention demonstration tests whether the knowledge layer can assign
a new candidate label to a high-$\Psi_s$ anomaly. The narrow result is software
capability: novelty detection and candidate schema recording. It is not, by
itself, proof of a new physical regime.

\subsection{Discovery D: vacuum hysteresis}
\texttt{experiments/notebooks/discovery\_vacuum\_hysteresis.py} tests whether a
high-flux pulse leaves localized $M_s$ memory. The selected output records a
center $M_s$ peak of 15.5843. The proper status is candidate simulation result
and measurement-protocol proposal.

\subsection{Discovery E: n=7 knot stability}
\texttt{experiments/notebooks/discovery\_knot\_symmetry.py} embeds a T(2,7)
constraint pattern and tracks a stability index toward the CODATA 2022
$\chi^*=137.035999177$ target. The release-selection note says to use it as a
weak higher-knot self-stabilization stress test, not a proof.

\subsection{Discovery F: adaptive-surface instability}
\texttt{experiments/notebooks/discovery\_adaptive\_surface.py} produces an
overloaded run. Because $\Psi_s$ grows from approximately 5.15 to approximately
979 in the selected output, this is presented as an instability boundary
and a safety test, not as an application success.

\section{Causal Graph from Run History}
\texttt{engine/causal\_graph.py} builds Granger-style directed edges from
simulation time series without an LLM. The CLI and optional web studio can
present these edges as triage hints: which series preceded which under the
current model. Optional LLM provider calls may later summarize a saved result,
but they do not create the edge, alter the finite audit, or turn a candidate
edge into validation. A persistent edge after decoupling controls remains an
artifact, not a discovery.

\section{Frontier and Gigamarathon Sweeps}
\texttt{experiments/outputs/frontier\_sweep.h5} summarizes 2,500 parameter
points and 2,487 non-trivial candidate records. This is hypothesis triage:
which regions deserve closer study. \texttt{gigamarathon\_1000\_results.h5}
covers 10 domains by 100 points each. This is coverage evidence for the
discovery engine, not physical validation.

\section{Promotion Ladder}
A candidate discovery moves through this ladder:
\begin{enumerate}
    \item simulation anomaly;
    \item reproducible runtime diagnostic;
    \item parameter-sensitivity report;
    \item falsification attempt;
    \item independent dataset or physical experiment;
    \item domain-specific validation.
\end{enumerate}

\section{Conclusion}
The discovery layer is valuable because it can search wider than a human can
manually search. Its danger is overstatement. The runtime papers present
discoveries as candidate simulation results with falsification handles until
external validation exists.
\begin{thebibliography}{9}
\bibitem{MujjabiI2026}
Steve Bico Mujjabi, MD. \emph{Vortex Stability in a Constrained Transport Vacuum and the Origin of the Fine-Structure Constant}. CDFD Part I Paper I, preprint, 2026.

\bibitem{MujjabiVI2026}
Steve Bico Mujjabi, MD. \emph{The Universal Torus Knot Hierarchy: $Z_n$ Symmetry, the Cinquefoil Family, and the General Structure of CDFT Particle Families}. CDFD Part I Paper VI, preprint, 2026.

\bibitem{MujjabiIX2026}
Steve Bico Mujjabi, MD. \emph{Even-$n$ Torus Knots in CDFT: The Exclusion of $T(2,2)$, the Extension to $T(2,2m)$, and the Anti-Phase Mass Scaling Law}. CDFD Part I Paper IX, preprint, 2026.

\bibitem{MujjabiX2026}
Steve Bico Mujjabi, MD. \emph{The Vacuum Equation of State: Deriving Torus Knot Self-Consistency from the Public CDFD Balance Law $\Psi = \Phi/C$}. CDFD Part I Paper X, preprint, 2026.

\bibitem{MujjabiXII2026}
Steve Bico Mujjabi, MD. \emph{Friction, Superconductivity, Plasma States, and Transport Mysteries: Observable Tests of Adaptive Constraint}. CDFD Part I Paper XII, preprint, 2026.

\bibitem{MujjabiPartII2026}
Steve Bico Mujjabi, MD. \emph{CDFD Part II: Origins of Life and Tri-Regime Bioenergetics}. Public release archive, 2026, \url{https://doi.org/10.5281/zenodo.20264779}.
\end{thebibliography}
\end{document}
